% =====================================================================
%  Chapter 32 — The Prior Art\idx{Prior art} and the Honest \idxw{Novelty} Claim
% =====================================================================

\chapter{The Prior Art and the Honest Novelty Claim}
\label{ch:prior-art}

A proposal that has been in the literature for twenty-three years cannot be
presented as new, and a book that pretended otherwise would be making a
claim its own chapters would contradict. This chapter maps the prior art\idx{Prior art}
and states, as precisely as the evidence allows, what is new here.

\section{The map}

Every identifier in this table was resolved against a registry before the
reference was used.

\begin{table}[!htbp]
\centering
\footnotesize
\caption{The prior-art map. Every candidate architecture in the original
concept, and its status.}
\label{tab:prior-art}
\begin{tabularx}{\textwidth}{@{}P{1.15in}P{1.55in}Y@{}}
\toprule
\rowcolor{panel}
\textbf{Candidate} & \textbf{Prior art, verified} & \textbf{Status} \\
\midrule
Lunar wide-field optical survey &
\code{arXiv:2611.05435}; CAPS NTRS 20030014831 and 20050186565 &
\textbf{Covered.} The lead-time table is published \\
\rowcolor{panel}
Lunar surface optical interferometer &
CAPS 2003, Table 1; \code{arXiv:2510.24901}; \code{arXiv:2408.04699} &
\textbf{Covered, and a mature, heavily funded field.} Not a gap \\
\rowcolor{panel}
Lunar bistatic radar &
\code{arXiv:1810.08712}; \code{arXiv:1908.05171} &
\textbf{Covered.} The range law is the limit, not the site \\
Cislunar optical survey and catalogue layer &
Lunar Sentinel, September 2026; \code{arXiv:2609.10663};
\code{arXiv:2604.10807} &
\textbf{Heavily covered} \\
\rowcolor{panel}
Space-based characterisation telescope &
\code{arXiv:2609.03371}; \code{arXiv:2609.02571}; NEO Surveyor; Roman &
\textbf{Heavily covered} \\
Low-Earth-orbit characterisation constellation &
\code{arXiv:2301.10348} &
\textbf{Covered} \\
\bottomrule
\end{tabularx}
\end{table}

\tabref{tab:prior-art} is the map, and it is not encouraging for the original
concept in its original form. Every architecture the owner's instinct
produced has been proposed, and most of them have been funded.

\section{What is genuinely not covered}

Three things are not in the literature, and they are of different kinds.

The first is negative, and it is the book's own correction. A lunar
\emph{thermal-infrared} station, with its background budget worked out, does
not appear in the literature. The 2026 extreme-environment study names
infrared monitoring stations and lunar polar observatories as asset classes
within a general Earth-Moon network and refers to the illumination conditions
and stable thermal environments near permanently shadowed regions, but it
specifies no surface mid-infrared instrument, no background budget, no power
architecture and no service model.

That gap is now closed, and closed against the lunar case: Chapter 23 works
the arithmetic out and the thermal channel loses to Sun-Earth L1. The
contribution here is the demonstration that the gap was a gap and that it
should not be filled.

The second is the substantive one. The lunar interferometry literature is
sized for \emph{astrometry}, and its 100-metre baselines are consistent
with that: the 2003 concept study's own design point is a visual magnitude
of 27 with a fringe visibility of 0.5, and its stated purpose is a drop in
the orbit-determination error. The recent lunar optical interferometry study
proposes a 100-metre baseline with 30-centimetre apertures over 450 to 850
nanometres, and the NIAC lunar imager is an astronomical instrument. The
2025 study spans ultraviolet to mid-infrared wavelengths but its baseline
design is optical, with a 30-centimetre aperture, a 100-metre baseline,
2-nanometre path measurement and a 5-metre main delay line.

Chapter 20's re-derivation establishes what the shape-imaging requirement
actually is, and it is not what those numbers are sized for. A
kilometre-scale optical baseline, matched to the target's angular size, for
the purpose of imaging rather than locating, does not appear in the lunar
literature.

The third is the one this book has derived and the literature does not
contain. The Earth-Moon baseline is a known quantity to millimetres. The
cislunar sensor-network modelling compares Earth-based, space-based and
lunar-surface-based observers and finds that lunar observers enjoy smaller
distances and better observability --- and names the Moon's horizon, the
necessity of clustering, and the logistics of operating a lunar station as
the drawbacks. That is an argument about observers.

It is not an argument about a \emph{ranging baseline}. The consequence of
putting the Earth and the Moon in the same frame and using the 384,400
kilometres between them as a parallax baseline is a range measurement with
no model in it, and the book found no published treatment of it for a
non-cooperative near-Earth target at close approach.

\section{What is claimed, precisely}

The claim is therefore bounded, and it is stated in the form the owner
asked for: \emph{a definite solution, novel if it can honestly be said to
be novel}.

\begin{table}[!htbp]
\centering
\small
\caption{What is claimed and what is not.}
\label{tab:claims}
\begin{tabularx}{\textwidth}{@{}P{2.57in}Y@{}}
\toprule
\rowcolor{panel}
\textbf{Claim} & \textbf{Status} \\
\midrule
The hardware is new & \textbf{No.} It is twenty-three years old and heavily
worked \\
\rowcolor{panel}
A lunar surface interferometer for asteroids & \textbf{No.} Proposed, and
proposed repeatedly \\
\rowcolor{panel}
A lunar thermal-infrared station with a worked budget & \textbf{No.} The
gap is demonstrated here and closed against the lunar case \\
\rowcolor{panel}
The shape-imaging baseline requirement, re-derived & \textbf{Yes.} The
concept study's own table reproduces as \(\lambda/(4\theta)\), and the
requirement follows \\
\rowcolor{panel}
The Earth-Moon two-station range to a non-cooperative target & \textbf{Yes,}
as far as this search found, and the derivation is published here \\
\rowcolor{panel}
The first quantitative argument for where in the planetary-defence
architecture a lunar platform earns its place & \textbf{Yes.} Together with
the error budgets and the selection rule that decide it \\
\bottomrule
\end{tabularx}
\end{table}

\tabref{tab:claims} is the honest statement, and its final row is the one
the book stands on. The contribution is of position and of quantification,
not of hardware. The analysis is new; the instrument is not.

\begin{seamnote}{Why a bounded claim is stronger than a broad one}
A broad claim invites the reader to check the hardware, where twenty-three
years of work will defeat it. A bounded claim invites the reader to check
the arithmetic, which is where the book has put everything it has. The
strongest defence of a proposal of this kind is that the re-derivation is
reproducible and the rule can be run against the reader's own targets and
declined where it should be.
\end{seamnote}

\begin{keytakeaways}
\item Every architecture in the original concept has been proposed, and most
have been funded; the original claim cannot be made.
\item The lunar interferometry literature is sized for astrometry, and the
kilometre-scale shape-imaging baseline requirement is not in it.
\item The Earth-Moon two-station range to a non-cooperative target is the
book's derived result, and no published treatment was found.
\item The claim is about position and quantification, not about hardware, and
it is stated in the form the evidence allows.
\end{keytakeaways}
